\section{Определение первообразной функции комплексного переменного.
Теорема о достаточных условиях существования первообразной. Теорема о
формуле Ньютона -Лейбница в случае комплексного
переменного.}\label{ux43eux43fux440ux435ux434ux435ux43bux435ux43dux438ux435-ux43fux435ux440ux432ux43eux43eux431ux440ux430ux437ux43dux43eux439-ux444ux443ux43dux43aux446ux438ux438-ux43aux43eux43cux43fux43bux435ux43aux441ux43dux43eux433ux43e-ux43fux435ux440ux435ux43cux435ux43dux43dux43eux433ux43e.-ux442ux435ux43eux440ux435ux43cux430-ux43e-ux434ux43eux441ux442ux430ux442ux43eux447ux43dux44bux445-ux443ux441ux43bux43eux432ux438ux44fux445-ux441ux443ux449ux435ux441ux442ux432ux43eux432ux430ux43dux438ux44f-ux43fux435ux440ux432ux43eux43eux431ux440ux430ux437ux43dux43eux439.-ux442ux435ux43eux440ux435ux43cux430-ux43e-ux444ux43eux440ux43cux443ux43bux435-ux43dux44cux44eux442ux43eux43dux430--ux43bux435ux439ux431ux43dux438ux446ux430-ux432-ux441ux43bux443ux447ux430ux435-ux43aux43eux43cux43fux43bux435ux43aux441ux43dux43eux433ux43e-ux43fux435ux440ux435ux43cux435ux43dux43dux43eux433ux43e.}

\textbf{Определение:} Функция \(F(z)\) называется \emph{первообразной
функции} \(f(z)\), определенной в области \(G\), если функция \(F(z)\)
определена в области \(G\), дифференцируема в этой области и \[
F'(z)=f(z),\quad z\in G
\]

\textbf{Теорема:} (Теорема о достаточных условиях существования
первообразной) Пусть

\begin{enumerate}
\def\labelenumi{\arabic{enumi}.}
\tightlist
\item
  \(f(z)\in C(G)\);
\item
  \(\forall\) замкнутого контура \(\gamma \in G\) выполняется условие
  \(\int\limits_{\gamma}f(z)\,dz=0\) Тогда \(f(z)\) имеет первообразную
  в \(G\).
\end{enumerate}

\textbf{Теорема:} (Теорема о формуле Ньютона-Лейбница в случае
комплексного переменного) Пусть

\begin{enumerate}
\def\labelenumi{\arabic{enumi}.}
\tightlist
\item
  \(f\in C(G), G\subset \mathbb{C}\)
\item
  \(\exists F(z)\) - первообразная функции \(f(z)\) в \(G\)
\item
  \(\gamma_{b,c}\) - любой путь (кусочно-гладкий), соединяющий точки
  \(b,c\in G\)
  (\(\gamma_{b,c}=\gamma_{b,c}(t), t\in[\alpha;\beta ], \ \gamma_{b,c}(t)\in C^{1}_{\text{кус.}}[\alpha;\beta], \ \gamma_{b,c}(\alpha)=b, \ \gamma_{b,c}(\beta)=c\))
  \[
  \int\limits_{\gamma_{b,c}}f(z)\,dz=F(c)-F(b),
  \]
\end{enumerate}
